\section{Migration, Versioning, and Governance}

Teams do not reach conformance through rewrite fantasy. They move by shrinking ambiguity. First add the audit in advisory mode. Then add root instructions, owner maps, test maps, generated zones, version manifests, and proof lanes. Then generate contracts, isolate data access, declare durable-truth ownership, and add rendered UX/security evidence for critical paths. Only after the repo can route and prove changes should teams accelerate agent-driven migration.

First-hour adoption must be no-write unless the operator explicitly applies a reviewed plan. \verb|jankurai adopt --mode observe| and \verb|jankurai init --dry-run| produce adoption or init artifacts without claiming conformance by accident. Migration analysis carries liability forward instead of hiding it: \verb|jankurai migrate --analyze| records missing tests, direct database risk, contract drift, and stack-distance evidence so a brownfield repo can improve by slices.

Governance needs a standard process, not private taste. Jankurai should use a JEP/RFC process for rule assignment, rule lifecycle, compatibility policy, badge policy, security disclosure, independent implementation certification, and old schema support. Rule lifecycle states are \texttt{draft}, \texttt{candidate}, \texttt{implementation-evidence}, \texttt{stable}, \texttt{superseded}, and \texttt{obsolete}. Draft rules need examples and fixtures. Candidate rules need implementation evidence. Stable rules need schema-backed output, documented repairs, and regression tests.

\begin{table}[t]
\caption{JEP/RFC governance gates.}
\label{tab:jep-governance}
\centering
\scriptsize
\begin{tabularx}{\columnwidth}{L{0.24\columnwidth} Y}
\toprule
\JKTableHeader
\textbf{Gate} & \textbf{Required evidence} \\
\midrule
Problem statement & Rule, schema, badge, or profile change names affected users and non-goals. \\
Compatibility review & Migration notes, schema impact, old-field interpretation, and downgrade behavior. \\
Implementation evidence & Reference implementation or independent implementation emits compatible artifacts. \\
Conformance fixtures & Positive and negative fixtures plus expected JSON. \\
Security/privacy review & Evidence classes, redaction, abuse risks, and disclosure path. \\
Release decision & Version bump, changelog, schemas, docs, paper provenance, and badge policy updated. \\
\bottomrule
\end{tabularx}
\end{table}

\begin{table}[t]
\caption{Schema compatibility policy.}
\label{tab:schema-compat}
\centering
\small
\begin{tabularx}{\columnwidth}{L{0.18\columnwidth} Y}
\toprule
\JKTableHeader
\textbf{Change} & \textbf{Meaning} \\
\midrule
Patch & Clarification, copy fix, example correction, or non-semantic schema note. \\
Minor & Additive fields or advisory rules that preserve existing report interpretation. \\
Major & Breaking report, witness, receipt, or conformance semantics. \\
\bottomrule
\end{tabularx}
\end{table}

\begin{table}[t]
\caption{Conformance claim tiers.}
\label{tab:claim-tiers}
\centering
\small
\begin{tabularx}{\columnwidth}{L{0.30\columnwidth} Y}
\toprule
\JKTableHeader
\textbf{Tier} & \textbf{Evidence} \\
\midrule
Self-attested & Repository emits local reports and witness artifacts. \\
CI-attested & Required CI uploads current receipts and merge witness for the commit. \\
Third-party-attested & Independent implementation or auditor reruns the suite and verifies artifacts. \\
Registry-listed & Public registry records version, level, evidence bundle, expiry, and compatibility. \\
\bottomrule
\end{tabularx}
\end{table}

Public evidence uses the same discipline. \texttt{registry} and \texttt{cell} expose reusable primitives; \texttt{bench}, \texttt{certify}, \texttt{govern}, and \texttt{publish} produce versioned benchmark, certification, governance, badge, and public-evidence artifacts. Those artifacts only matter because \texttt{agent/standard-version.toml} binds them to a standard version, auditor version, schema version, paper edition, and manifest profile label. \texttt{paper\_edition} is provenance for the paper, not a core conformance field.

\begin{table}[t]
\caption{Version taxonomy and release checklist.}
\label{tab:version-taxonomy}
\centering
\scriptsize
\begin{tabularx}{\columnwidth}{L{0.28\columnwidth} Y}
\toprule
\JKTableHeader
\textbf{Versioned item} & \textbf{Release rule} \\
\midrule
Standard version & Changes normative conformance semantics, levels, or rule obligations. \\
Schema version & Changes report, receipt, witness, manifest, or queue interpretation. \\
Rule registry version & Adds, stabilizes, supersedes, or retires HLT families. \\
Profile version & Changes non-normative stack or architecture guidance only. \\
Paper edition & Changes explanatory prose, evidence presentation, citations, or companion artifacts. \\
Release checklist & Run conformance, paper, score, schema tests, diff check, changelog, and compatibility notes. \\
\bottomrule
\end{tabularx}
\end{table}

Streaming migration follows the same governance. Do not begin by replacing Kafka. Begin by hiding event buses behind contracts and adapters, measuring broker readiness with replay and compatibility tests, and making the migration path visible in \texttt{agent/boundaries.toml}. Streaming clients belong behind adapters and generated event contracts; the broker choice is profile evidence, not Jankurai Core.

\begin{table}[t]
\caption{CI adoption modes.}
\label{tab:ci-modes}
\centering
\small
\begin{tabularx}{\columnwidth}{L{0.20\columnwidth} Y}
\toprule
\JKTableHeader
\textbf{Mode} & \textbf{Merge behavior} \\
\midrule
Advisory & Emit JSON/Markdown reports and never block. \\
Guarded & Block critical caps and malformed output. \\
Standard & Block high and critical findings plus required receipts. \\
Ratchet & Prevent score or finding regression without a dated waiver. \\
Release & Gate shipped artifacts on audit, tests, security, contracts, DB, e2e, versions, and release evidence. \\
\bottomrule
\end{tabularx}
\end{table}

Governance should separate empirical claims from policy choices. Scoring weights, hard caps, conformance levels, RPNs, reference-profile scores, and rule IDs are versioned policy. Public source claims require citations. New audit output fields require schema versioning. Breaking compliance changes require migration notes and example repairs. The goal is not to freeze one 2026 opinion forever; it is to make the standard itself updatable without letting every repository invent a private dialect.

Independent implementations should be welcome and testable. A non-Rust auditor may claim compatibility when it emits schema-valid artifacts, uses the closed decision enum, preserves receipt invalidation semantics, passes the conformance fixtures for the claimed version, and documents any optional extension fields without changing Core interpretation.
